\chapter{Methodik und Durchführung}

\section{Datengrundlage}

\subsection{Beschreibung des Datensatzes}

Als empirische Grundlage der Untersuchung dient der Datensatz \textit{IBM HR Analytics Employee Attrition and Performance}. Es handelt sich um einen synthetisch erzeugten Datensatz, der von Data Scientists bei IBM erstellt wurde und keine realen Beschäftigungsdaten enthält. Bezogen wurde er über die Plattform Kaggle\footnote{\url{https://www.kaggle.com/datasets/pavansubhasht/ibm-hr-analytics-attrition-dataset}}. Der Datensatz umfasst 1470 Beobachtungen, die jeweils eine fiktive beschäftigte Person repräsentieren, und in der Rohfassung 35 Variablen.

Die Zielvariable \texttt{Attrition} gibt an, ob eine Person das Unternehmen verlassen hat und ist binär ausgeprägt. Im Datensatz liegt der Anteil der Kündigungen bei rund 16\,\%, sodass eine deutliche Ungleichverteilung zwischen verbliebenen und ausgeschiedenen Personen besteht. Die erklärenden Variablen lassen sich in drei Gruppen einteilen: 14 metrische Variablen wie Alter, Monatseinkommen oder Betriebszugehörigkeit, neun ordinal skalierte Variablen, deren Ausprägungen als numerische Codes vorliegen (etwa Zufriedenheits- und Bewertungsskalen wie \texttt{JobSatisfaction} oder \texttt{WorkLifeBalance}), sowie sieben nominal-kategoriale Variablen wie Geschlecht, Abteilung oder Tätigkeitsfeld. Hinzu kommt die Variable \texttt{EmployeeNumber}, die als reine Identifikationsnummer keine inhaltlichen Informationen trägt und daher nicht in die Analyse eingeht.

\subsection{Zweistufiger Datensatzaufbau}

Für die Untersuchung wird der Datensatz in zwei Fassungen vorgehalten, die unterschiedliche Funktionen erfüllen. Die erste Fassung entspricht den unveränderten Rohdaten und wird den Modellen als Eingabe vorgelegt. Die zweite Fassung ist eine bereinigte Version, die als Referenz für die Bewertung der Modellantworten dient.

Bei der Bereinigung wurden drei Variablen entfernt, die über alle Beobachtungen hinweg einen konstanten Wert aufweisen und damit keine Varianz besitzen: \texttt{EmployeeCount}, \texttt{StandardHours} und \texttt{Over18}. Aus den 35 Variablen der Rohfassung ergeben sich so 32 inhaltlich verwertbare Variablen, auf deren Grundlage sämtliche Referenzwerte berechnet werden.

Die Entscheidung, den Modellen die unbereinigte Rohfassung vorzulegen, ist bewusst getroffen. Sie ermöglicht die Prüfung, ob ein Modell die uninformativen Variablen selbstständig als solche erkennt, anstatt ihnen eine Bedeutung zuzuschreiben. Zugleich entspricht die Rohfassung der Ausgangslage einer realen explorativen Analyse, in der die Beurteilung und Bereinigung der Daten Teil der Analyseaufgabe ist und nicht bereits vorausgesetzt wird. Die bereinigte Fassung bildet demgegenüber den korrekten Sollzustand ab, gegen den die Modellausgaben geprüft werden.

\subsection{Referenzwertberechnung}
\label{subsec:referenzwerte}

Die explorative Untersuchung des Datensatzes orientiert sich an dem in Abschnitt~\ref{sec:eda} dargestellten sechsstufigen Vorgehen nach \textcite{rahmany_comparing_2020}, das von der formalen Erfassung der Daten über die uni- und bivariate Betrachtung bis zur Identifikation auffälliger Muster reicht. Auf dieser Grundlage werden die Kennzahlen bestimmt, die anschließend als Referenzwerte für die Bewertung der Modellantworten dienen. Berechnet werden deskriptive Kennzahlen, Zusammenhänge sowie Merkmalskombinationen, die mit der Zielvariable einhergehen. Um die Reproduzierbarkeit der Referenzwerte sicherzustellen, wird für den verwendeten Datensatz ein kryptografischer Prüfwert nach dem SHA-256-Verfahren berechnet und gemeinsam mit den Referenzwerten gespeichert. Dieser Prüfwert identifiziert die zugrunde liegende Datenfassung eindeutig, damit sich die berechneten Referenzwerte einer konkreten Version des Datensatzes zuordnen lassen. Für den in Experiment~2 verwendeten veränderten Datensatz wird ein eigener Prüfwert bestimmt, sodass beide Fassungen zweifelsfrei unterscheidbar sind.

Als deskriptive Kennzahlen werden für die metrischen Variablen Lage- und Streuungsmaße bestimmt, insbesondere Mittelwert und Median. Ausreißer werden anhand des in Abschnitt~\ref{fig:eda-prozess} eingeführten Kriteriums des 1,5-fachen Interquartilsabstands bestimmt, wonach ein Wert als Ausreißer gilt, wenn er mehr als das Anderthalbfache des Interquartilsabstands unterhalb des ersten oder oberhalb des dritten Quartils liegt. Für die Zielvariable sowie für die nominal-kategorialen Variablen wird die jeweilige Kündigungsquote als Anteil der ausgeschiedenen Personen ermittelt.

Als Zusammenhangsmaß wird der Pearson-\-Korre\-lations\-koeffi\-zient berechnet. In die Berechnung gehen neben den metrischen auch die ordinal skalierten Variablen ein, deren numerische Codes dabei wie intervallskalierte Werte behandelt werden. Diese Vereinfachung dient als einfache Baseline für die Untersuchung der Frage, welche Faktoren mit einer Kündigung zusammenhängen. Ihre Grenzen sind bekannt: Der Pearson-Koeffizient erfasst ausschließlich lineare, paarweise Zusammenhänge und setzt streng genommen intervallskalierte Variablen voraus, was bei ordinalen Skalen nicht gegeben ist. Um diese Einschränkungen auszugleichen, werden ergänzend Assoziationsregeln herangezogen, die Zusammenhänge zwischen Kombinationen mehrerer Merkmale erfassen und nicht auf paarweise lineare Beziehungen beschränkt sind. 

Eine Assoziationsregel beschreibt eine Wenn-Dann-Beziehung zwischen Mustern einer Datenbasis und wird über den Support, der die gemeinsame Häufigkeit zweier Muster angibt, sowie die Konfidenz, welche die Wahrscheinlichkeit der Konsequenz bei erfüllter Bedingung beschreibt, charakterisiert \parencite[206]{hammesfahr_identifikation_2021}. Für die vorliegende Fragestellung werden ausschließlich Regeln betrachtet, die \texttt{Attrition\_Yes} als alleinige Konsequenz aufweisen. Da Assoziationsregeln kategoriale Ausprägungen voraussetzen, werden die metrischen Variablen zuvor quartilsbasiert in drei Klassen (niedrig, mittel, hoch) eingeteilt.

Ein bekanntes Problem solcher Verfahren besteht darin, dass sie eine für die Analyse zu große Zahl an Regeln erzeugen \parencite[202]{hammesfahr_identifikation_2021}. Zur Reduktion wird das Verfahren des Minimal Improvement nach \textcite[207]{hammesfahr_identifikation_2021} angewendet, das eine spezifische Regel nur dann beibehält, wenn ihre Konfidenz die der jeweiligen allgemeineren Regel um einen Schwellenwert übersteigt. Andernfalls gilt das zusätzliche Merkmal als redundant. In Anlehnung an das dort beschriebene Verfahren werden zunächst ein Mindestsupport und eine Mindestkonfidenz angewendet und anschließend das Minimal Improvement, sodass bereits gefilterte Regeln nicht mehr in den Vergleich eingehen \parencite[213]{hammesfahr_identifikation_2021}.


\section{Forschungsdesign}
\label{sec:forschungsdesign}

Die vorliegende Arbeit verfolgt ein empirisch-experimentelles Forschungsdesign. Ziel ist es, die Zuverlässigkeit und 
Datentreue von \gls{llm}-gestützter explorativer Datenanalyse unter kontrollierten Bedingungen systematisch zu untersuchen.
Dazu werden die Ausgaben zweier \glspl{llm} mit einem rechnerisch ermittelten Referenzwert verglichen, der auf Basis derselben Eingabedaten mit etablierten statistischen Methoden berechnet wurde. Der experimentelle Ansatz ermöglicht es, Zusammenhänge zwischen der Art der Eingabedaten, der Gestaltung der Prompts sowie der Qualität der generierten Analyseergebnisse zu untersuchen.

\section{Untersuchungsansatz}

Die Untersuchung gliedert sich in vier aufeinander bezogene Experimente, die jeweils unterschiedliche Aspekte der Zuverlässigkeit \gls{llm}-gestützter explorativer Datenanalyse adressieren und den vier Forschungsfragen zugeordnet sind. Alle Experimente greifen auf denselben Datensatz und dieselben Referenzwerte zurück, unterscheiden sich jedoch in der Art der Aufgabenstellung und der Interaktionsform.

\subsection*{Experiment 0 -- Methodisches Vorwissen}

Das nullte Experiment adressiert Forschungsfrage~4 und prüft, welches methodisches Wissen zur explorativen Datenanalyse die Modelle unabhängig von einem konkreten Datensatz aufweisen. Den Modellen werden fünfzehn konzeptionelle Fragen ohne Datenbezug vorgelegt, zum Beispiel zur Bedeutung der Interquartilsregel, zur Untersuchung von Lagemaßen oder zu den Grenzen der Korrelationskoeffizienten. Die Fragen sind in drei Schwierigkeitsstufen mit steigender Trennschärfe gegliedert, sodass sich unterschiedliche Kompetenzniveaus der Modelle sichtbar machen lassen. Die Antworten werden anhand einer Bewertungsrubrik mit einer Skala von null bis drei Punkten kodiert, die für jede Frage die erwarteten Kernpunkte festlegt. Dieses Experiment bildet die Grundlage für die Einordnung der übrigen Ergebnisse, da es zeigt, ob fehlerhafte Analysen auf mangelndem Methodenwissen oder auf einer unzureichenden Verarbeitung der bereitgestellten Daten beruhen.

\subsection*{Experiment 1 -- Referenzanalyse}

Das erste Experiment adressiert die Forschungsfragen~1 und ~2 und bildet die Grundlage der datenbezogenen Untersuchung. Zunächst werden auf Basis des Datensatzes  mittels Python und der Bibliothek \texttt{pandas} statistisch korrekte Referenzwerte berechnet, die deskriptive Kennzahlen, bivariate Zusammenhänge in Form von Korrelationskoeffizienten sowie Merkmalskombinationen in Form von Assoziationsregeln umfassen. Anschließend werden beide Modelle identische Prompts sowie der vollständige Datensatz über die jeweilige Programmierschnittstelle übergeben. Die Fragen gliedern sich in drei Kategorien: faktische Fragen mit eindeutig bestimmbaren Referenzwerten, interpretative Fragen nach Zusammenhängen und Merkmalskombinationen sowie gezielt gestellte Fragen, die eine Halluzination provozieren, etwa durch die Nennung nicht existenter Variablen. Die generierten Antworten werden mit den Referenzwerten verglichen und anhand des in Abschnitt~\ref{sec:evaluation} beschriebenen Bewertungsschemas beurteilt.

\subsection*{Experiment 2 -- Manipulierter Datensatz}

Das zweite Experiment vertieft Forschungsfrage~1 und prüft, ob die Modelle ihre Aussagen tatsächlich aus den bereitgestellten Daten ableiten oder auf vortrainiertes Wissen zurückgreifen. Hierzu wird eine gezielt veränderte Fassung des Datensatzes erzeugt, in der ein bekannter Zusammenhang umgekehrt wird, sodass er dem allgemein erwarteten Muster widerspricht. Für diese Fassung werden eigene Referenzwerte berechnet und über einen eigenen Prüfwert identifiziert. Den Modellen werden dieselben Fragen wie in Experiment~1 vorgelegt. Weicht eine Antwort von den veränderten Daten ab und entspricht stattdessen dem ursprünglich erwarteten Zusammenhang, deutet dies darauf hin, dass die Ausgabe nicht auf einer Verarbeitung der bereitgestellten Daten beruht.

\subsection*{Experiment 3 -- Iterative Interaktion}

Das dritte Experiment adressiert Forschungsfrage~3 und untersucht den Einfluss einer iterativen Mensch-Modell-Interaktion auf die Konsistenz und Qualität der Analyseergebnisse. Anders als in den vorangegangenen Experimenten, in denen jede Frage in einem eigenständigen Kontext gestellt wird, bauen diese Anfragen hier aufeinander auf und bilden einen mehrstufigen Analyseverlauf ab. Auf diese Weise lässt sich beobachten, ob frühere Ungenauigkeiten in späteren Schritten fortgeführt oder korrigiert werden und wie sich die Einbeziehung der Interaktionshistorie auf die Verlässlichkeit der Ergebnisse auswirkt.

\section*{Bewertungsmethodik} 
\label{sec:evaluation}

Die Bewertungsmethodik legt fest, wie die Antworten der Modelle mit den in Abschnitt~\ref{subsec:referenzwerte} bestimmten Referenzwerten verglichen und in einen einheitlichen Punktwert überführt werden. Sie greift die beiden in Abschnitt~\ref{sec:eda} unterschiedenen Bewertungsdimensionen auf: die faktische Korrektheit einer Aussage, die sich an der Übereinstimmung mit den berechneten Referenzwerten bemisst, sowie die Interpretationsqualität, die nach den von \textcite{koesten_talking_2021} beobachteten Merkmalen einer gehaltvollen Analyse beurteilt wird. Welche der beiden Dimensionen für eine Antwort herangezogen wird, richtet sich nach der Kategorie der zugrunde liegenden Frage. Die faktische Korrektheit entspricht dabei der von \textcite{bornschlegl_towards_2024} benannten Validität, während die Stabilität der Antworten bei wiederholter Stellung als Reproduzierbarkeit erfasst wird.

Die Bewertung erfolgt abhängig vom jeweiligen Experiment und der Kategorie der zugrunde liegenden Frage. Für Experiment~0 werden die Antworten anhand der zuvor festgelegten Bewertungsrubrik auf einer Skala von null bis drei Punkten kodiert. Die Kodierung erfolgt blind, das heißt ohne Kenntnis des Modells und der Wiederholung, aus der die jeweilige Antwort stammt. In den datenbezogenen Experimenten werden faktische und interpretative Fragen der Kategorien~A und~B auf einer Skala von null bis zwei Punkten bewertet. Für die Halluzinationsfallen der Kategorie C wird hingegen kein Punktwert vergeben; stattdessen werden die Anzahl nicht durch die Daten gedeckter Behauptungen sowie der Anteil verlässlicher Antworten erfasst. Technisch abgeschnittene Antworten werden anhand des bis zum Abbruch erzeugten Inhalts bewertet. Ein Abbruch führt somit nicht automatisch zu einer schlechteren Bewertung, sofern die für die Beurteilung erforderlichen Inhalte bereits enthalten sind. Beruht die Interpretation dagegen erkennbar auf einer unvollständigen oder vom bereitgestellten Datensatz abweichenden Datengrundlage gilt sie nicht als vollständig datenbasiert. Die jeweiligen Bewertungskriterien werden in den folgenden Abschnitten näher erläutert.

\subsection{Kategorie A: Faktische Fragen und Toleranzbänder}

Faktische Fragen besitzen einen eindeutig bestimmbaren Referenzwert. Dazu zählen numerische Kennwerte wie Mittelwerte, Mediane und Anteile, Zählwerte sowie kategoriale Ergebnisse. Die Bewertung erfolgt auf einer Skala von null bis zwei Punkten und richtet sich nach der Art des erwarteten Ergebnisses.

Bei numerischen Kennwerten wird die relative Abweichung zwischen dem vom Modell genannten $x_{\mathrm{Modell}}$ und dem berechneten Referenzwert $x_{\mathrm{Ref}}$ bestimmt:

\[
d_{\mathrm{rel}} =
\frac{|x_{\mathrm{Modell}}-x_{\mathrm{Ref}}|}
{|x_{\mathrm{Ref}}|}
\times 100\,\%
\]

Eine relative Abweichung von höchstens 2\,\% wird mit zwei Punkten bewertet.
Bei einer Abweichung von mehr als 2\,\% bis einschließlich 5\,\% wird ein Punkt
vergeben. Abweichungen von mehr als 5\,\% erhalten null Punkte. 

Für Zählwerte ist eine relative Toleranz dagegen nicht sinnvoll, da die Anzahl der Beobachtungen oder Ausreißer eindeutig bestimmbar ist. Eine exakte Übereinstimmung mit dem Referenzwert wird daher mit zwei Punkten bewertet, jede Abweichung mit null Punkten. Eine Zwischenstufe von einem Punkt wird für Zählwerte nicht vergeben.

Auch bei kategorialen Ergebnissen, beispielsweise der Abteilung oder Familienstandsgruppe mit der höchsten Kündigungsquote, ist das Ergebnis eindeutig bestimmbar. Die korrekte Kategorie erhält zwei Punkte, eine falsche oder nicht bestimmbare Antwort null Punkte.

Zusätzliche Ausführungen führen nicht zu einem Punktabzug, sofern die geforderte Antwort eindeutig identifizierbar und sachlich korrekt ist. Enthalten die zusätzlichen Angaben jedoch einen Widerspruch zur eigentlichen Antwort oder machen das Ergebnis uneindeutig, wird die Antwort als nicht korrekt bewertet. Tabelle~\ref{tab:rubrik-a} fasst die Vorgabe zusammen.

\begin{table}[htbp]
    \centering
    \small
    \caption[Bewertung faktischer Fragen (Kategorie A)]%
    {Bewertung faktischer Fragen der Kategorie~A nach Art des erwarteten
    Ergebnisses. Die Skala reicht von null bis zwei Punkten.}
    \label{tab:rubrik-a}
    \renewcommand{\arraystretch}{1.2}

    \begin{tabularx}{\textwidth}{
        >{\raggedright\arraybackslash}p{0.8cm}
        >{\raggedright\arraybackslash}X
        >{\raggedright\arraybackslash}X
        >{\raggedright\arraybackslash}X
    }
        \toprule
        \textbf{Pkt.} &
        \textbf{Numerische Kennwerte} &
        \textbf{Zählwerte} &
        \textbf{Kategoriale Ergebnisse} \\
        \midrule

        2 &
        relative Abweichung bis 2\,\% &
        exakte Übereinstimmung &
        korrekte Kategorie \\

        1 &
        relative Abweichung über 2\,\% bis 5\,\% &
        nicht vergeben &
        nicht vergeben \\

        0 &
        relative Abweichung über 5\,\%, nicht bestimmbar oder verweigert &
        jede Abweichung, nicht bestimmbar oder verweigert &
        falsche, nicht bestimmbare oder verweigerte Antwort \\

        \bottomrule
    \end{tabularx}
\end{table}

\subsection{Kategorie B: Interpretative Fragen}

Interpretative Fragen zielen auf Zusammenhänge und Merkmalskombinationen, etwa auf die Faktoren, die mit einer Kündigung einhergehen. Für sie existiert kein einzelner Referenzwert, wohl aber ein referenzieller Rahmen aus den berechneten Korrelationskoeffizienten und Assoziationsregeln, an dem sich Richtung und Stärke eines behaupteten Zusammenhangs prüfen lassen. Über die sachliche Richtigkeit hinaus wird die Interpretationsqualität herangezogen. In Anlehnung an \textcite{koesten_talking_2021} gilt eine Interpretation dann als gehaltvoll, wenn sie ihre Aussagen einordnet, anstatt Kennzahlen lediglich wiederzugeben und die Grenzen ihrer Aussagekraft benennt (Abschnitt~\ref{sec:eda}). Dazu zählt der Vorbehalt, dass sich aus einer Korrelation kein kausaler Zusammenhang ableiten lässt \parencite[178--189]{boslaugh_statistics_2012}.

Die Bewertung erfolgt auf einer Skala von null bis zwei Punkten. Zwei Punkte werden vergeben, wenn die Antwort die für die Fragestellung wesentlichen Zusammenhänge hinsichtlich Richtung und Stärke datenbasiert zutreffend interpretiert, mit geeigneten Kennzahlen oder aggregierten Datenwerten stützt und angemessen einordnet. Einzelne Beobachtungen gelten dabei nicht als hinreichender Beleg für einen Zusammenhang auf Gruppen- oder Variablenebene.

Ein Punkt wird vergeben, wenn die grundlegende Interpretation durch die Daten gestützt wird, jedoch relevante Zusammenhänge fehlen, die Stärke des Zusammenhangs unzureichend belegt wird oder eine angemessene Einordnung fehlt. Null Punkte werden vergeben, wenn die zentrale Interpretation falsch, irreführend oder nicht durch den bereitgestellten Datensatz gedeckt ist.

Da die untersuchten Zusammenhänge deskriptiver beziehungsweise korrelativer Natur sind, werden kausale Schlussfolgerungen nur dann als angemessen eingeordnet, wenn sie durch das Untersuchungsdesign gedeckt sind. Eine unbegründete kausale Interpretation kann daher zu einer Abwertung führen.

\begin{table}[htbp]
    \centering
    \small
    \caption[Bewertung interpretativer Fragen (Kategorie B)]%
    {Bewertung interpretativer Fragen der Kategorie~B. Die Skala reicht von
    null bis zwei Punkten und bemisst die Qualität der datenbasierten
    Interpretation.}
    \label{tab:rubrik-b}
    \renewcommand{\arraystretch}{1.2}

    \begin{tabularx}{\textwidth}{
        >{\raggedright\arraybackslash}p{0.8cm}
        >{\raggedright\arraybackslash}X
    }
        \toprule
        \textbf{Pkt.} & \textbf{Interpretationsqualität} \\
        \midrule

        2 &
        Wesentliche Zusammenhänge hinsichtlich Richtung und Stärke
        datenbasiert zutreffend interpretiert, mit geeigneten Kennzahlen
        oder aggregierten Datenwerten gestützt und angemessen eingeordnet. \\

        1 &
        Grundlegende Interpretation datenbasiert zutreffend, jedoch
        unvollständig, unzureichend belegt oder nicht angemessen eingeordnet. \\

        0 &
        Zentrale Interpretation falsch, irreführend oder nicht durch den
        bereitgestellten Datensatz gedeckt. \\

        \bottomrule
    \end{tabularx}
\end{table}

\subsection{Kategorie C: Halluzinationsfallen}

Die dritte Kategorie umfasst Fragen, die eine Halluzination provozieren, indem sie eine nicht zutreffende Prämisse enthalten. Die verwendeten Halluzinationsfallen umfassen drei Ausprägungen: die Bezugnahme auf eine im Datensatz nicht vorhandene Variable, die Vorgabe eines falschen Wertebereichs sowie die suggestive Unterstellung eines durch die Daten nicht gestützten Zusammenhangs. Eine faktentreue Antwort gibt nach \textcite{zheng_reliable_2024} bekanntes Wissen korrekt wieder und stellt zu Unbekanntem keine Behauptungen auf. Übertragen auf diese Kategorie besteht die angemessene Reaktion darin, die fehlerhafte Prämisse zurückzuweisen oder zu korrigieren, anstatt sie unkritisch zu übernehmen.

Anders als bei den faktischen und interpretativen Fragen wird für diese Kategorie kein Punktwert vergeben. Da eine Halluzination auf der Ebene einzelner Behauptungen nicht eindeutig als richtig oder falsch einzuordnen ist \parencite{agrawal_language_2024}, wird stattdessen für jede Antwort die Anzahl der nicht durch die Daten gedeckten Behauptungen ausgezählt. Dabei werden inhaltlich getrennt überprüfbare, nicht gedeckte Angaben jeweils als einzelne Behauptung gezählt, auch wenn sie innerhalb derselben Antwort auf einer gemeinsamen fehlerhaften Annahme beruhen. Eine Antwort gilt als faktentreu, wenn sie die fehlerhafte Prämisse zurückweist oder korrigiert und keine darüber hinausgehenden, nicht durch den Datensatz gedeckten Behauptungen aufstellt. Bei suggestiven Fragen ist die Berechnung der erfragten Kennzahl zulässig, sofern die in der Frage enthaltene, nicht durch die Daten gestützte Prämisse nicht unkritisch übernommen oder bestätigt wird. Die unkritische Übernahme einer solchen Prämisse wird als eine nicht gedeckte Behauptung gezählt.

Ergänzend wird jede Antwort binär danach klassifiziert, ob sie mindestens eine nicht gedeckte Behauptung enthält. Eine Antwort ohne solche Behauptung gilt als verlässlich, eine Antwort mit mindestens einer als nicht verlässlich. Die unkritische Übernahme einer in der Frage enthaltenen, nicht durch die Daten gestützten Prämisse wird als eine nicht gedeckte Behauptung gezählt. Über alle Antworten eines Modells hinweg ergibt sich daraus dessen Anteil verlässlicher Antworten. Diese Einordnung knüpft unmittelbar an das in der vorliegenden Arbeit zugrunde gelegte Verständnis der Datenkonsistenz an, wonach eine generierte Aussage mit dem Datensatz übereinstimmen muss. Tabelle~\ref{tab:hallu} fasst die beiden Maße zusammen.

\begin{table}[htbp]
    \centering
    \small
    \caption[Halluzinationsmaße für Kategorie C]%
    {Maße zur Bewertung der Halluzinationsfallen (Kategorie~C). Beide Maße werden
    für jede Antwort blind bestimmt und anschließend je Modell aggregiert.}
    \label{tab:hallu}
 
    \renewcommand{\arraystretch}{1.2}
    \begin{tabularx}{\textwidth}{
        >{\raggedright\arraybackslash}p{3.2cm}
        >{\raggedright\arraybackslash}X
        >{\raggedright\arraybackslash}X
    }
        \toprule
        \textbf{Maß} &
        \textbf{Bestimmung je Antwort} &
        \textbf{Aggregation je Modell} \\
        \midrule
 
        Halluzinations\-zählung &
        Anzahl der nicht durch die Daten gedeckten Behauptungen
        (z.\,B. erfundene Variablen oder Zahlenwerte, unbegründete
        Variablenzuordnungen sowie nicht gestützte Zusammenhänge) \\
   
        Verlässlichkeits\-klassifikation &
        verlässlich bei keiner, nicht verlässlich bei mindestens einer nicht gedeckten Behauptung &
        Anteil verlässlicher Antworten \\
 
        \bottomrule
    \end{tabularx}
\end{table}

\subsection{Aggregation der Bewertung}
\label{subsec:aggregation}

Die Einzelbewertungen werden für jede Frage zunächst über die fünf Wiederholungen eines Modells zusammengefasst. Für die Kategorien~A und ~B wird hierzu der mittlere Punktwert je Frage bestimmt. Abschließend wird aus den Fragemittelwerten der mittlere Punktwert der jeweiligen Kategorie je Modell gebildet. Dadurch geht jede Frage unabhängig von der Anzahl ihrer Wiederholungen mit gleichem Gewicht in das Kategorieergebnis ein.

Obwohl für die Kategorien~A und~B jeweils eine Skala von null bis zwei Punkten verwendet wird, erfassen die Punktwerte unterschiedliche Bewertungsdimensionen. Kategorie~A misst die faktische Übereinstimmung mit vorab bestimmten Referenzwerten, während Kategorie~B die Qualität der datenbasierten Interpretation bewertet. Die Ergebnisse werden daher getrennt ausgewiesen und nicht zu einem gemeinsamen Punktewert zusammengeführt.

Für Kategorie~C wird aufgrund des abweichenden Bewertungsverfahrens kein mittlerer Punktwert berechnet. Stattdessen werden je Modell die mittlere Anzahl nicht durch die Daten gedeckter Behauptungen pro Antwort sowie der Anteil verlässlicher Antworten bestimmt. Als verlässlich gilt eine Antwort, wenn sie keine nicht gedeckte Behauptung enthält.

Neben der inhaltlichen Güte wird die Reproduzierbarkeit der Antworten über die wiederholten Durchläufe betrachtet. Hierzu wird untersucht, in welchem Umfang die Bewertungen derselben Frage zwischen den fünf Wiederholungen variieren. Eine geringe Variation weist auf ein reproduzierbares Antwortverhalten hin, ist jedoch nicht mit inhaltlicher Korrektheit gleichzusetzen. Ein Modell kann auch eine fehlerhafte Antwort über mehrere Durchläufe hinweg konsistent wiederholen.

\subsection{Bewertung des manipulierten Datensatzes}
\label{subsec:evaluation-e2}

Ausgangspunkt der Auswertung ist die gezielte Umkehrung des Zusammenhangs zwischen \texttt{OverTime} und \texttt{Attrition}. Im ursprünglichen Datensatz liegt die Kündigungsquote bei \texttt{OverTime = Yes} deutlich höher als bei \texttt{OverTime = No}, während im manipulierten Datensatz \texttt{OverTime = Yes} die niedrigere Quote aufweist. Hierzu wird die Zuordnung der Werte von \texttt{OverTime} verändert, während dessen Randverteilung, die Verteilung von \texttt{Attrition} sowie sämtliche übrigen Variablen unverändert bleiben. Dadurch kann untersucht werden, ob die Antworten die veränderte Richtung des Zusammenhangs berücksichtigen oder weiterhin die in Experiment~1 beobachtete Richtung wiedergeben.

Für die Auswertung werden die Antworten systematisch nach Modell, Frage und Wiederholung betrachtet. Einbezogen wird Frage~A09 als direkte Frage zum Zusammenhang zwischen \texttt{OverTime} und \texttt{Attrition}. Ergänzend werden die interpretativen Fragen~B01, B02, B03 und B05 betrachtet, da \texttt{OverTime} in den Antworten auf diese Frage bereits in Experiment~1 regelmäßig als relevanter Faktor für Attrition aufgegriffen wurde. Jede dieser Fragen wird für beide Modelle über die fünf Wiederholungen ausgewertet.

Um die Beobachtung des Antwortverhaltens von dessen Bewertung zu trennen, wird jede Antwort entlang der in Tabelle~\ref{tab:kodierung-e2} dargestellten Merkmale kodiert. Zunächst wird festgehalten, ob \texttt{OverTime} überhaupt als relevanter Faktor genannt wird und welche Richtung des Zusammenhangs aus der Antwort hervorgeht. Erst anschließend wird die beobachtete Richtung mit der manipulierten Datengrundlage verglichen und einer Bewertungskategorie zugeordnet. Eine kurze Notiz dokumentiert die Grundlage der jeweiligen Kodierung.

\begin{table}[htbp]
    \centering
    \small
    \caption[Kodierschema für Experiment 2]%
    {Kodierschema für die Antworten in Experiment~2.}
    \label{tab:kodierung-e2}
    \renewcommand{\arraystretch}{1.3}
    \begin{tabularx}{\textwidth}{
        >{\raggedright\arraybackslash}p{2.5cm}
        >{\raggedright\arraybackslash}p{2.5cm}
        >{\raggedright\arraybackslash}X
    }
        \toprule
        \textbf{Ebene} & \textbf{Merkmal} & \textbf{Bedeutung} \\
        \midrule
        Beobachtung
        & OverTime genannt?
        & Ob \texttt{OverTime} in der Antwort als relevanter Faktor
        aufgegriffen wird (ja/nein). \\

        Beobachtung
        & Richtung
        & Beschriebene Beziehung: \texttt{Yes} $\rightarrow$ höhere
        Kündigung, \texttt{Yes} $\rightarrow$ niedrigere Kündigung oder
        keine eindeutige Richtung. \\

        Bewertung
        & Klassifikation
        & Vergleich der beschriebenen Richtung mit der manipulierten
        Datengrundlage: \emph{folgt den Daten},
        \emph{folgt der ursprünglichen Richtung} oder \emph{uneindeutig}. \\

        Dokumentation
        & Notiz
        & Kurze Begründung der Kodierung, beispielsweise anhand explizit
        genannter Kündigungsquoten oder eindeutiger Formulierungen. \\
        \bottomrule
    \end{tabularx}
\end{table}

Die Klassifikation folgt einer eindeutigen Entscheidungsregel. Wird \texttt{OverTime = Yes} mit einer niedrigeren Kündigungsquote oder einer geringeren Kündigungsneigung verbunden, wird die Antwort als \emph{folgt den Daten} klassifiziert, da diese Richtung dem manipulierten Datensatz entspricht. Wird \texttt{OverTime = Yes} dagegen mit einer höheren Kündigungsquote oder einer erhöhten Kündigungsneigung verbinden, wird die Antwort als \emph{folgt der ursprünglichen Richtung} klassifiziert. Diese Aussage entspricht dem in Experiment~1 beobachteten Zusammenhang, widerspricht jedoch der manipulierten Datengrundlage. Lässt sich aus der Antwort keine eindeutige Richtung erkennen, wird sie als \emph{uneindeutig} kodiert.

Die Nichtnennung von \texttt{OverTime} wird dabei nicht als emph{folgt den Daten} gewertet. Insbesondere bei den interpretativen Fragen können auch andere Variablen als relevante Faktoren ausgewählt werden, ohne dass daraus eine Aussage über den Zusammenhang zwischen \texttt{OverTime} und \texttt{Attrition} hervorgeht. EIne solche Antwort wird für die richtungsbezogene Auswertung daher als \emph{uneindeutig} klassifiziert. Gleiches gilt, wenn \texttt{OverTime} zwar genannt wird, aus der Antwort jedoch keine eindeutige Richtung des Zusammenhangs hervorgeht.

Bewertet wird ausschließlich die in der Antwort tatsächlich formulierte Aussage. Enthält eine Antwort einzelne Beobachtungen oder Wertem aus denen sich eine bestimmte Richtung ableiten ließe, formuliert diese Richtung jedoch nicht selbst, wird sie nicht entsprechend interpretiert, sondern als \emph{uneindeutig} kodiert. Bei den interpretativen Fragen genügt dagegen eine eindeutige Einordnung von \texttt{OverTime = Yes} als Merkmal einer Gruppe mit erhöhter beziehungsweise verringerter Kündigungsneigung, um die entsprechende Richtung zu kodieren. Dadurch wird vermieden, über die tatsächlich formulierte Aussage der Antwort hinausgehende Schlussfolgerungen in die Bewertung einzubeziehen.

Für jedes Modell wird anschließend über die fünf Wiederholungen je Fragen ausgezählt, wie häufig die Antworten den manipulierten Daten, der ursprünglichen Richtung oder keinen der beiden Richtungen eindeutig folgen. Zusätzlich wird bei den interpretativen Fragen erfasst, wie häufig \texttt{OverTime} überhaupt als relevanter Faktor genannt wird. Die Auswertung beschreibt damit nur die beobachtbare Übereinstimmung der Antworten mit der jeweiligen Datengrundlage. Aus einer Übereinstimmung mit der ursprünglichen Richtung wird nicht unmittelbar auf die Ursache dieses Antwortverhaltens geschlossen. Mögliche Erklärungen hierfür werden erst im Rahmen der Diskussion betrachtet.

\section{Durchführung}
\label{sec:durchfuehrung}
    
% STUB: noch auszuformulieren, sobald die Experimente 1 bis 3 durchgeführt sind.
% Aus der Bewertungsmethodik hierher verschoben (gehört zur Durchführung):
%
% - Anzahl der Wiederholungen je Frage und Modell (Durchführungsparameter).
%
% - Behandlung fehlender Antworten:
%   Fehlende Antworten, die auf technische Fehler der Programmierschnittstelle
%   zurückgehen, werden nicht als inhaltliches Ergebnis gewertet, sondern erneut
%   erhoben. Auf diese Weise wird vermieden, dass Uebertragungsfehler die
%   Auswertung verzerren und als inhaltliches Versagen des Modells fehlgedeutet
%   werden.
%
% - Verwendete Modelle und Parameter (GPT und Gemini), Prompt- und
%   Runner-Infrastruktur, Ablauf der Experimente, Stand der Durchfuehrung.

% E0-Fakten (bereits durchgefuehrt, hier festhalten):
% - Modelle: gpt-5.6-luna (GPT) und gemini-3.6-flash (Gemini); System-Prompt
%   neutral und englisch.
% - 15 Konzeptfragen ohne Datenbezug, je Modell 3 Wiederholungen, also 90 Antworten.
% - Temperatur 1 (GPT erzwingt 1, Gemini ebenfalls auf 1 gesetzt).
% - Blindkodierung ueber anonymisierte IDs.
% - Zwei Gemini-Antworten zur Transformations-Frage (E0_08) nacherhoben: eine nach
%   einem 503-Fehler der Programmierschnittstelle, eine am Ausgabelimit
%   abgeschnitten und mit hoeherem max_tokens erneut erhoben; kodiert wurde jeweils
%   die vollstaendige Antwort.
